The top-level $\CauchyFilterSubbasisUp$ route consumes the finite $\FilterBaseUp$ generation rows of \autoref{ch:concrete-instances-filterbase-namecert} before it reads the $\CauchyFilterUp$ packet of \autoref{ch:concrete-instances-cauchyfilter-namecert}. Only after that base-to-Cauchy handoff does the route pass through the $\StreamNameUp$ finite window, $\RegSeqRatUp$ readback, and terminal $\RealUp$ seal of \autoref{ch:concrete-instances-streamname-namecert}, \autoref{ch:concrete-instances-regseqrat-namecert}, and \autoref{ch:concrete-instances-real-namecert}. This is a finite generation route, not a quotient filter, selected limit, countable-choice schedule, or ambient $\RealUp$ completeness claim.

\begin{theorem}[CauchyFilterSubbasis Cauchy-completion sibling route]
\label{thm:cauchy-filter-subbasis-cauchy-completion-sibling-route}
Every accepted $\CauchyFilterSubbasisUp$ packet exported to a completion-facing consumer first reads the Cauchy-completion criterion surface of \autoref{ch:concrete-instances-cauchy-filter-completion-criterion-namecert}. The required dependency is concrete: subbasis neighborhoods are generated from the finite $\FilterBaseUp$ rows, the Cauchy filter completion criterion supplies the completion rows, $\StreamNameUp$ and $\RegSeqRatUp$ provide only finite windows and rational readback, and the terminal $\RealUp$ seal records nonescape from selected limits, quotient filters, countable choice, and ambient completeness.
\end{theorem}
\begin{proof}
Project the top-level route described above. The subbasis neighborhood row is read before any Cauchy packet is exported, so the consumer begins with finite $\FilterBaseUp$ generation and the displayed $\CauchyFilterUp$ handoff.

Next read the Cauchy-completion criterion row. It supplies the completion-facing criterion used by the sibling chapter, while the $\StreamNameUp$ and $\RegSeqRatUp$ rows restrict every read to finite stream windows and rational approximants.

The final $\RealUp$ seal is only a nonescape row. It does not select a limit, identify quotient filters, introduce countable choice, or assert ambient completeness beyond the finite criterion route already displayed.
\end{proof}
